%------------------------------------------------------------------------------%
\documentclass[fleqn,utf8,aspectratio=169,12pt]{beamer}

\usepackage{gaal}

\title{Posições Relativas entre Planos e Retas}

%------------------------------------------------------------------------------%
\begin{document}

\addtitle

%------------------------------------------------------------------------------%
\begin{frame}{Posições Relativas entre uma Reta e um Plano}
  \onslide<+->{Existem três possibilidades}
  \begin{itemize}[<+->]
      \item a reta está contida no plano
      \item a reta é paralela ao plano e não está contida nele
      \item a reta é secante ao plano
  \end{itemize}
\end{frame}

%------------------------------------------------------------------------------%
\begin{frame}{Posições Relativas entre uma Reta e um Plano}
  Considere a reta
  \tabto{40mm}
  \(
    r\colon \vec{r} = \vec{p} + t\vec{v}
  \)

  \pause
  \medskip
  e o plano
  \tabto{40mm}
  \(
    \pi\colon ax + by + cz + d = 0
  \)

  \pause
  \smallskip
  com vetor normal
  \tabto{40mm}
  \(
    \vec{n} = (a, b, c)
  \)

  \pause
  \medskip
  Comparamos o vetor diretor $\vec{v}$ da reta
  \pause
  com o vetor normal $\vec{n}$ do plano
\end{frame}

%------------------------------------------------------------------------------%
\begin{frame}{Critério Usando o Vetor Diretor}
\begin{columns}[t]
\column{0.3\textwidth}
  Se
  \[
    \vec{v}\cdot\vec{n} \neq 0
  \]
  \pause
  a reta é \emph{secante} ao plano
\column{0.5\textwidth}
\pause
  Se
  \[
    \vec{v}\cdot\vec{n} = 0
  \]
  \pause
  a reta é paralela ou está contida no plano

  \bigskip
  \bigskip
  \pause
  Se houver um ponto em comum a reta está contida no plano
\end{columns}
\end{frame}

%------------------------------------------------------------------------------%
%------------------------------------------------------------------------------%
\begin{question}
  Determine a posição relativa entre a reta e o plano \\
  e o ponto de interseção, se ele existir
  \[
    r:
    \begin{cases}
      x = 1 +  t \\
      y = 2 -  t \\
      z = 3 + 2t
    \end{cases}
  \]
  \[
    \gamma: x + y + z - 6 = 0
  \]
\end{question}

%------------------------------------------------------------------------------%
\begin{resolution}{Posição Relativa}
  \onslide<+->{O vetor diretor da reta é \qquad
  \(
    \vec{v} = (1, -1, 2)
  \)}

  \onslide<+->{O vetor normal do plano é \quad
  \(
    \vec{n} = (1, 1, 1)
  \)}
  \begin{align*}
    \onslide<+->{\vec{v} \cdot \vec{n}}
    \onslide<+->{& = 1 \times 1 + (- 1)1 + 2 \times 1} \\
    \onslide<+->{& = 1 - 1 + 2} \\
    \onslide<+->{& = 2} \\
    \onslide<+->{& \neq 0}
  \end{align*}
  \onslide<+->{a reta é secante ao plano}
\end{resolution}

%------------------------------------------------------------------------------%
\begin{resolution}{Ponto de Interseção}
\begin{columns}
\column{0.35\textwidth}
  \onslide<+->{Substituir a equação da reta
  \[
    \begin{cases}
      x = 1 +  t \\
      y = 2 -  t \\
      z = 3 + 2t
    \end{cases}
  \]}
  \onslide<+->{na equação do plano}
\column{0.5\textwidth}
  \begin{align*}
    \onslide<+->{x + y + z - 6              & = 0 \\}
    \onslide<+->{(1+t) + (2-t) + (3+2t) - 6 & = 0 \\}
    \onslide<+->{t - t + 2t + 1 + 2 + 3 - 6 & = 0 \\}
    \onslide<+->{2t                         & = 0 \\}
    \onslide<+->{t                          & = 0}
  \end{align*}
\end{columns}
\end{resolution}

%------------------------------------------------------------------------------%
\begin{resolution}{Ponto de Interseção}
  \onslide<+->{Coordenadas do ponto com $t=0$}
  \begin{align*}
    \onslide<+->{x & = 1 + t }
    \onslide<+->{  \;=\; 1 \\}
    \onslide<+->{y & = 2 - t}
    \onslide<+->{  \;=\; 2 \\}
    \onslide<+->{z & = 3 + 2t}
    \onslide<+->{  \;=\; 3}
  \end{align*}

  \onslide<+->{\[
    P = (1, 2, 3)
  \]}
\end{resolution}

%------------------------------------------------------------------------------%
\endinput

%------------------------------------------------------------------------------%
\begin{question}
  Determine a posição relativa entre a reta e o plano \\
  e o ponto de interseção, se ele existir
  \[
    r\colon
    \begin{cases}
      x = 1 +  t \\
      y = 2 + 2t \\
      z = 1 -  t
    \end{cases}
  \]
  \[
    \gamma\colon x + y - z + 10 = 0
  \]
\end{question}

%------------------------------------------------------------------------------%
\begin{resolution}{Posição Relativa}
  \onslide<+->{O vetor diretor da reta é \qquad
  \(
    \vec{v} = (1, 2, -1)
  \)}

  \onslide<+->{O vetor normal do plano é \quad
  \(
    \vec{n} = (1, 1, -1)
  \)}
  \begin{align*}
    \onslide<+->{\vec{v} \cdot \vec{n}}
    \onslide<+->{& = 1 \times 1 + 2\times 1 + (-1)(-1)} \\
    \onslide<+->{& = 1 + 2 - 1} \\
    \onslide<+->{& = 2} \\
    \onslide<+->{& \neq 0}
  \end{align*}
  \onslide<+->{a reta é secante ao plano}
\end{resolution}

%------------------------------------------------------------------------------%
\begin{resolution}{Ponto de Interseção}
\begin{columns}
\column{0.35\textwidth}
  \onslide<+->{Substituir a equação da reta
  \[
    \begin{cases}
      x = 1 +  t \\
      y = 2 + 2t \\
      z = 1 -  t
    \end{cases}
  \]}
  \onslide<+->{na equação do plano}
\column{0.5\textwidth}
  \begin{align*}
    \onslide<+->{x + y - z + 10              & = 0 \\}
    \onslide<+->{(1+t) + (2+2t) - (1-t) + 10 & = 0 \\}
    \onslide<+->{t + 2t + t + 1 + 2 - 1 + 10 & = 0 \\}
    \onslide<+->{4t + 12                     & = 0 \\}
    \onslide<+->{t                           & = \frac{12}{4} \\}
    \onslide<+->{                            & = 3}
  \end{align*}
\end{columns}
\end{resolution}

%------------------------------------------------------------------------------%
\begin{resolution}{Ponto de Interseção}
  \onslide<+->{Coordenadas do ponto com $t=0$}
  \begin{align*}
    \onslide<+->{x & = 1 + t}
    \onslide<+->{  \;=\; 1 + 3}
    \onslide<+->{  \;=\; 4 \\}
    \onslide<+->{y & = 2 + 2t}
    \onslide<+->{  \;=\; 2 + 2\times 3}
    \onslide<+->{  \;=\; 8 \\}
    \onslide<+->{z & = 1 - t}
    \onslide<+->{  \;=\; 1 - 3}
    \onslide<+->{  \;=\; -2}
  \end{align*}

  \onslide<+->{\[
    P = (4, 8, -2)
  \]}
\end{resolution}

%------------------------------------------------------------------------------%
\endinput


%------------------------------------------------------------------------------%
\begin{question}
  Determine a posição relativa entre a reta e o plano \\
  e o ponto de interseção, se ele existir
  \[
    r\colon
    \begin{cases}
      x = 1 + t \\
      y = 2 - t \\
      z = 3
    \end{cases}
  \]
  \[
    \gamma\colon x + y - 3 = 0
  \]

\end{question}

%------------------------------------------------------------------------------%
\begin{resolution}{Posição Relativa}
  \onslide<+->{O vetor diretor da reta é \qquad
  \(
    \vec{v} = (1, -1, 0)
  \)}

  \onslide<+->{O vetor normal do plano é \quad
  \(
    \vec{n} = (1, 1, 0)
  \)}
  \begin{align*}
    \onslide<+->{\vec{v} \cdot \vec{n}}
    \onslide<+->{& = 1 \times 1 + (- 1)1 + 0 \times 0} \\
    \onslide<+->{& = 1 - 1 + 0} \\
    \onslide<+->{& = 0}
  \end{align*}
  \onslide<+->{a reta é paralela ao plano ou está contida nele}
\end{resolution}

%------------------------------------------------------------------------------%
\begin{resolution}{Ponto de Interseção}
  Escolhemos o ponto da reta
  \[
    A = (1, 2, 3)
  \]
  \pause
  e substituimos na equação do plano
  \pause
  \[
    x + y - 3
    \pause
    = 1 + 2 - 3
    \pause
    = 0
  \]
  \pause
  Como o ponto da reta satisfaz a equação do plano
  \[
    \pause
    A\in\gamma
  \]
\end{resolution}

%------------------------------------------------------------------------------%
\begin{resolution}{Ponto de Interseção}
  Como a reta possui um ponto no plano e é paralela ao plano

  \pause
  todos os seus pontos pertencem ao plano
  \pause
  \[
    r\subset\gamma
  \]
  A reta está \emph{contida no plano}
\end{resolution}

%------------------------------------------------------------------------------%
\endinput


%%------------------------------------------------------------------------------%
%% \framedgraphic{Gráfico de uma Função Real}{B_01_grafico_funcao_real}{}
%
%%------------------------------------------------------------------------------%
%\section{Lista Mínima}
%
%%------------------------------------------------------------------------------%
%\begin{frame}{Lista Mínima}
%  \begin{enumerate}
%    \item
%  \end{enumerate}
%
%  \vfill
%  Atenção: A prova é baseada no livro, não nas apresentações
%\end{frame}

%------------------------------------------------------------------------------%
\end{document}
%------------------------------------------------------------------------------%
